\section{本章小结}
\label{sec:opt-lp-summary}

回到章首的作坊。等值线在可行四边形上平移、最后停在一个角点上，这幅画面在本章变成了一套完整的计算。标准型把各种线性模型统一成$\bm A\bm x=\bm b$、$\bm x\ge\bm 0$；角点是基本可行解，由选定$m$列解一个线性方程组得到；最优值有限时，最优解可以在这有限个基本可行解中找到。单纯形法用约化费用判断哪条边能改进，用最小比值检验决定走多远，从一个基换到相邻的基；两阶段法提供起点，Bland规则排除循环。停止时留下的乘子$\bm y$就是原料的价格：它与最优解同值，证明了最优；第一阶段的乘子在不可行时组合出矛盾；一个可行点加一条下降射线证明无下界。三种结局各有一份不依赖求解过程、只需矩阵乘法就能核对的证书。

本章反复区分了几个容易混淆的结论，它们所需的条件各不相同，表~\ref{tab:opt-lp-summary}把它们放在一起。可行域有界只是最优解存在的充分条件，不是必要条件；线性规划特有的是，只要可行且有下界，最优值就能取到，标准型中还能取在顶点上。有限终止依赖防循环规则和精确算术，它不蕴含多项式时间；数值求解停止、残差很小与严格认证最优是三个不同层次的判断。

\begin{table*}[!htbp]
  \centering
  \caption{本章主要结论及其条件。“标准型”指$\min\bm c^\top\bm x$，$\bm A\bm x=\bm b$，$\bm x\ge\bm 0$；最右列为正文中的位置。}
  \label{tab:opt-lp-summary}
  \begin{tabular}{@{}lll@{}}
    \toprule
    结论 & 所需条件 & 位置 \\
    \midrule
    最优值能取到，且有最优基本可行解 & 标准型可行，目标有下界 & 定理~\ref{thm:opt-lp-optimal-bfs} \\
    基本可行解即极点 & $\bm A$满行秩 & 定理~\ref{thm:opt-lp-bfs-extreme} \\
    约化费用非负推出全局最优 & 当前基可行 & 定理~\ref{thm:opt-lp-optimality-criterion} \\
    单纯形法有限终止 & Bland规则（或其他防循环规则），精确算术 & 定理~\ref{thm:opt-lp-finite} \\
    强对偶，两侧最优解都存在 & 原问题可行且有下界；不需要内点条件 & 定理~\ref{thm:opt-lp-strong-duality} \\
    不可行证书$\bm A^\top\bm z\ge\bm 0$，$\bm b^\top\bm z<0$ & 原约束系统无解 & 定理~\ref{thm:opt-lp-farkas} \\
    无下界证书$(\bm x_0,\bm d)$ & 已知可行点，$\bm A\bm d=\bm 0$，$\bm d\ge\bm 0$，$\bm c^\top\bm d<0$ & 第~\ref{sec:opt-lp-certificates}节 \\
    影子价格给出最优值的变化率 & 右端变化后原基本解仍非负 & 第~\ref{sec:opt-lp-rhs}节 \\
    凸包上线性最优值不变 & 候选集有限，目标线性，凸包精确 & 定理~\ref{thm:opt-lp-hull-value} \\
    \bottomrule
  \end{tabular}
\end{table*}

得到最优解之后的两个问题改变的是不同的对象。参数敏感性仍在同一类连续模型中，回答数据改变时原基还能用多久：右端变化时检查原基本解是否非负，最优值的变化率是影子价格；目标系数变化时检查约化费用是否非负，变化率是原始解。两种情况都只在各自的范围内有效，越过折点就要换基。线性松弛则把另一个离散问题放进连续多面体：精确凸包保持线性最优值，外逼近只给出一个界，恢复原候选还需要依据原约束单独设计，并重新检查可行性。

单纯形法沿可行域的边界从顶点走到顶点。另一条路线是从可行域内部出发，沿一条光滑的路径逼近最优点，这就是内点法。它的每一步需要求解牛顿方程，并借助障碍函数保持在内部，还要处理可行起点与停止依据；这些内容依赖光滑优化的工具，放在第~\ref{chap:constrained-optimization}章集中展开。本章的求解与认证已经由单纯形法和对偶证书完整闭合。
